\section{System Architecture}

The Automated Shelf Monitoring System consists of three main components: image capture, computer vision processing, and alert/dashboard presentation.

\subsection{Component Overview}

\begin{itemize}
    \item \textbf{Image Capture Module:} Receives shelf images from a fixed camera source.
    \item \textbf{Vision Processing Module:} Applies object detection and shelf occupancy analysis to each image.
    \item \textbf{Alert and Dashboard Module:} Presents occupancy results, active alerts, and history to store staff.
\end{itemize}

\subsection{Data Flow}

\begin{enumerate}
    \item The camera captures a shelf image at a configured interval.
    \item The image is forwarded to the vision processing module.
    \item The detection model identifies products and determines occupied shelf regions.
    \item Occupancy results are stored in the status database.
    \item If a region becomes empty, an alert is generated and displayed in the dashboard.
\end{enumerate}

\begin{figure}[H]
    \centering
    \includegraphics[width=0.7\textwidth]{shelfsense_system_architecture.png}
    \caption{System Architecture Diagram}
    \label{fig:system-architecture}
\end{figure}

\section{Design Decisions}

\subsection{Model Selection}

A lightweight object detection model was chosen to support near-real-time processing. The system uses YOLOv8n (nano) - optimized for speed - with transfer learning from COCO pre-trained weights. The model was fine-tuned on a custom dataset of retail SKUs.
\\
YOLOv8 was selected because:
\begin{itemize}
    \item It provides real-time performance with acceptable accuracy on edge devices.
    \item It offers multiple model variants (n/s/m) for trade-off between speed and accuracy.
    \item It has excellent support for deployment on various platforms including edge devices.
\end{itemize}

\subsection{Shelf Region Representation}

The shelf is represented as a set of monitoring regions corresponding to product segments or slots. This simplifies the occupancy problem and makes the alerting logic easy to interpret for store staff.

\subsection{Alert Logic}

Alerts are generated when a previously occupied region is detected as empty in the latest processed image. To reduce false alarms, the system uses a short temporal smoothing window, requiring the region to remain empty for two consecutive captures before raising an alert.

\section{Module Descriptions}

\subsection{Image Capture Module}

The image capture module accepts:
\begin{itemize}
    \item Periodic frames from a live camera feed, or
    \item Manual captures of a live frame at any instant.
\end{itemize}

Captured images are normalized and resized before analysis.

\subsection{Vision Processing Module}

The vision processing module performs the following tasks:
\begin{itemize}
    \item Preprocessing: resize, normalize, and crop input images.
    \item Product detection: locate and classify products on the shelf.
    \item Region occupancy classification: determine whether each shelf region is occupied.
    \item Result aggregation: create a structured status object for the dashboard.
\end{itemize}

\begin{figure}[H]
    \centering
    \includegraphics[width=0.8\textwidth]{shelfsense_detection_pipeline.png}
    \caption{Detection Pipeline}
    \label{fig:detection-pipeline}
\end{figure}

\subsection{Alert and Dashboard Module}

This module maintains current status and alert history. It provides:
\begin{itemize}
    \item A shelf occupancy view with colored indicators for occupied and empty regions.
    \item A list of active alerts with timestamp, region identifier, and recommended action.
    \item A history of resolved alerts.
\end{itemize}

\section{Data Design}

The system stores the following entities:
\begin{itemize}
    \item \textbf{Image Record:} capture timestamp, image path, capture source.
    \item \textbf{Shelf Region Status:} region ID, occupancy state, product SKU, confidence score.
    \item \textbf{Alert:} alert ID, region ID, alert type, created timestamp, resolved status.
\end{itemize}

\section{Technical Stack}

The prototype uses a Python-based vision backend with a lightweight detection model, and a simple dashboard interface for visualization. The backend can be deployed on a laptop or edge device for pilot retail use.
